\chapter{Patterns of Recursion and Dictionaries}

This chapter consolidates the recursive patterns seen so far and applies them
to a very common abstract data type: the \textbf{dictionary}.

\section{A catalogue of patterns}
Most list functions written so far fall into one of a few shapes. Recognizing
the shape tells you most of the code before you write it.

\begin{center}
\begin{tabular}{lll}
\toprule
Pattern & Base case & Recursive case \\
\midrule
Combine (sum, count) & identity value & \code{(op (f (first l)) (recur (rest l)))} \\
Search (contains?) & \code{false} & \code{(or (test (first l)) (recur (rest l)))} \\
All-satisfy & \code{true} & \code{(and (test (first l)) (recur (rest l)))} \\
Transform (map) & \code{empty} & \code{(cons (f (first l)) (recur (rest l)))} \\
Select (filter) & \code{empty} & \code{cons} first only if it passes the test \\
Accumulate & accumulator & recur on rest with updated accumulator \\
Find-first & \code{false} & first if it passes, else recur \\
\bottomrule
\end{tabular}
\end{center}

These shapes are exactly what higher-order functions (\code{foldr},
\code{ormap}, \code{andmap}, \code{map}, \code{filter}, \code{foldl}) capture in
Chapter~\ref{ch:abstraction}.

\section{Dictionaries}

\begin{definition}[Dictionary]
A \textbf{dictionary} (also map, table, associative array) stores
\emph{key--value pairs} where each key occurs at most once. Operations:
\begin{itemize}
  \item \textbf{lookup}: given a key, produce its value (or indicate absence);
  \item \textbf{add}: add a key and value (or replace the value of an existing key);
  \item \textbf{remove}: delete the pair with a given key.
\end{itemize}
\end{definition}
Examples: student ID $\to$ name, word $\to$ definition, symbol $\to$ count.

\begin{center}
\begin{tikzpicture}[k/.style={font=\small},arr/.style={-{Stealth},thick,blue!50!black}]
  \draw[fill=blue!5] (0,0) ellipse (1.2 and 1.5);
  \draw[fill=green!6] (5,0) ellipse (1.2 and 1.5);
  \node[font=\scriptsize\bfseries] at (0,1.8) {keys (unique)};
  \node[font=\scriptsize\bfseries] at (5,1.8) {values};
  \node[k] (a) at (0,0.8) {12345};
  \node[k] (b) at (0,0) {31415};
  \node[k] (c) at (0,-0.8) {27182};
  \node[k] (x) at (5,0.8) {"Ada"};
  \node[k] (y) at (5,0) {"Grace"};
  \node[k] (z) at (5,-0.8) {"Alan"};
  \draw[arr] (a) -- (x); \draw[arr] (b) -- (y); \draw[arr] (c) -- (z);
\end{tikzpicture}
\end{center}
The dictionary is an \textbf{abstract data type} (ADT): it is defined by its
operations, and may be implemented in different ways with different
efficiency.

\section{Association lists}

The simplest implementation is a list of two-element lists.

\begin{lstlisting}[style=racket]
;; An Association List (AL) is one of:
;; * empty
;; * (cons (list Num Str) AL)
;;   requires: each key (Num) is unique

(define marks '((12345 "Ada") (31415 "Grace") (27182 "Alan")))
\end{lstlisting}

In memory, \code{marks} is a list whose elements are themselves two-element
lists. A lookup walks down the outer list comparing keys:

\begin{center}
\begin{tikzpicture}[c/.style={draw,minimum width=7mm,minimum height=6mm,inner sep=1pt},
  w/.style={draw,minimum height=6mm,inner sep=2pt,font=\small\ttfamily},
  ptr/.style={-{Stealth}},
  nil/.style={path picture={\draw (path picture bounding box.south west) -- (path picture bounding box.north east);}}]
  \foreach \i/\k/\v in {0/12345/"Ada",1/31415/"Grace",2/27182/"Alan"} {
    \node[c] (f\i) at (\i*3.6,0) {};
    \fill (f\i.center) circle (1.2pt);
    \node[w] (k\i) at (\i*3.6+0.35,-1.3) {\k};
    \node[c] (kr\i) at ($(k\i.east)+(0.35,0)$) {};
    \node[w] (v\i) at ($(kr\i.east)+(0.75,-0.9)$) {\v};
    \node[c,nil] (vr\i) at ($(v\i.east)+(0.35,0)$) {};
    \draw[ptr] (f\i.center) -- (f\i.center |- k\i.north);
    \draw[ptr] (kr\i.center) |- (v\i.west);
  }
  \foreach \i in {0,1,2} \node[c] (r\i) at (\i*3.6+0.7,0) {};
  \node[c,nil] at (2*3.6+0.7,0) {};
  \draw[ptr] (r0.center) -- (f1.west);
  \draw[ptr] (r1.center) -- (f2.west);
  \node[font=\scriptsize\ttfamily,anchor=east] at (-0.5,0) {marks};
  \node[font=\scriptsize,text=red!60!black] at (0.9,-0.65) {$\ne 31415$};
  \node[font=\scriptsize,text=green!40!black] at (4.5,-0.65) {$= 31415$};
\end{tikzpicture}
\end{center}

A slash in a rest box stands for \code{empty}.

\paragraph{Template.}
\begin{lstlisting}[style=racket]
;; al-template: AL -> Any
(define (al-template alst)
  (cond [(empty? alst) ...]
        [else (... (first (first alst)) ...     ; key
                   (second (first alst)) ...    ; value
                   (al-template (rest alst)) ...)]))
\end{lstlisting}

\subsection{Lookup}
\begin{lstlisting}[style=racket]
;; (lookup-al k alst) produces the value associated with k in alst,
;;   or false if k is not a key.
;; Examples:
(check-expect (lookup-al 31415 marks) "Grace")
(check-expect (lookup-al 99999 marks) false)

;; lookup-al: Num AL -> (anyof Str false)
(define (lookup-al k alst)
  (cond [(empty? alst) false]
        [(= k (first (first alst))) (second (first alst))]
        [else (lookup-al k (rest alst))]))
\end{lstlisting}
Producing \code{false} for ``not found'' is a common CS 135 convention; the
contract uses \code{(anyof Str false)} to document it.

\paragraph{Trace.}
\begin{lstlisting}[style=racket]
(lookup-al 3 (list (list 1 "a") (list 3 "c")))
=> (cond [(empty? (list (list 1 "a") (list 3 "c"))) false] ...)
...  ; first clause false, second: (= 3 1) is false
=> (lookup-al 3 (rest (list (list 1 "a") (list 3 "c"))))
=> (lookup-al 3 (list (list 3 "c")))
...  ; (= 3 3) is true
=> (second (first (list (list 3 "c"))))
=> (second (list 3 "c"))
=> "c"
\end{lstlisting}

\subsection{Add or update}
\begin{lstlisting}[style=racket]
;; (add-al k v alst) produces an AL like alst but with k mapped to v,
;;   replacing any previous value for k.
;; Examples:
(check-expect (add-al 1 "x" empty) '((1 "x")))
(check-expect (add-al 1 "y" '((1 "x") (2 "z"))) '((1 "y") (2 "z")))
(check-expect (add-al 3 "w" '((1 "x"))) '((1 "x") (3 "w")))

;; add-al: Num Str AL -> AL
(define (add-al k v alst)
  (cond [(empty? alst) (list (list k v))]
        [(= k (first (first alst))) (cons (list k v) (rest alst))]
        [else (cons (first alst) (add-al k v (rest alst)))]))
\end{lstlisting}

\subsection{Remove}
\begin{lstlisting}[style=racket]
;; (remove-al k alst) produces alst without the pair with key k.
;; remove-al: Num AL -> AL
(check-expect (remove-al 2 '((1 "a") (2 "b") (3 "c"))) '((1 "a") (3 "c")))
(check-expect (remove-al 9 '((1 "a"))) '((1 "a")))
(define (remove-al k alst)
  (cond [(empty? alst) empty]
        [(= k (first (first alst))) (rest alst)]
        [else (cons (first alst) (remove-al k (rest alst)))]))
\end{lstlisting}

\subsection{Efficiency}
Each operation may traverse the whole list, so all three take time
proportional to $n$, the number of pairs. Keeping the list sorted by key does
not help asymptotically for a linked list (we still have to walk to the
position). Binary search trees (Chapter~\ref{ch:trees}) bring lookup down to
proportional to the tree's height.

\section{Application: word frequencies}
\begin{lstlisting}[style=racket]
;; A FreqTable is a (listof (list Sym Nat)) with unique symbols.

;; (count-words los) produces a FreqTable giving how often each symbol
;;   occurs in los.
;; count-words: (listof Sym) -> FreqTable
(check-expect (count-words '(a b a)) '((a 2) (b 1)))
(define (count-words los)
  (cond [(empty? los) empty]
        [else (increment (first los) (count-words (rest los)))]))

;; (increment s ft) adds 1 to the count of s in ft (adding s with
;;   count 1 if absent), keeping s at the front.
;; increment: Sym FreqTable -> FreqTable
(define (increment s ft)
  (cond [(empty? ft) (list (list s 1))]
        [(symbol=? s (first (first ft)))
         (cons (list s (add1 (second (first ft)))) (rest ft))]
        [else (cons (first ft) (increment s (rest ft)))]))
\end{lstlisting}
(Order of the result depends on processing order; tests should match the
order your design specifies, or the purpose should promise none.)

\section{Pairs: lists or structures?}
A key--value pair could be \code{(list k v)} or \code{(make-entry k v)}.
Lists are convenient with list abbreviations and quoting; structures give
named selectors and a type predicate. CS 135 uses both; be consistent within a
program and say which in the data definition.

\section{Summary}
\begin{itemize}
  \item Many list functions follow a small number of recursive patterns.
  \item A dictionary maps unique keys to values with lookup, add, remove.
  \item Association lists implement dictionaries in linear time per operation.
  \item ``Not found'' is commonly signalled by producing \code{false}.
\end{itemize}
